\subsection{One notation for donors, learner fields and contexts}\label{m:states}
We distinguish a saved runner state from the trajectory subsequently generated from it. For training seed $i$, $P_i$ is the complete shared onset runner at update $w_0=390$; $C_i^{+}$ is the unimpaired fork at removal $w_1=440$; and $D_i$ is the impaired runner at that same removal update, with corruption switched off. The superscript $+$ reminds the reader that $C_i^{+}$ contains 50 additional clean updates that are \emph{not} available in $P_i$.

Write a composed removal state as $\mathcal S(a,v,o;x)$: $a$ is actor parameters, $v$ critic parameters, $o$ the Adam state including its internal counter, and $x$ every other numerical runner field. In particular, $x_D$ retains the natural impaired environment, observations, future training random key, rollout-update counter, TrainState step and auxiliary fields. Keeping TrainState step is not the same as keeping Adam's internal count: the latter follows the selected $o$ donor. Figure~\ref{m:counterfactuals} gives the complete construction used by the primary contrasts.

\begin{align}
N_i &= \operatorname{Cont}\{\mathcal S(a_D,v_D,o_D;x_D)\}, &
R_i &= \operatorname{Cont}\{\mathcal S(a_P,v_D,o_D;x_D)\},\\
Q_{D,i} &= \operatorname{Cont}\{\mathcal S(a_P,v_P,o_P;x_D)\}, &
B_{D,i} &= \operatorname{Cont}\{\mathcal S(a_D,v_P,o_P;x_D)\}.
\end{align}
Here $\operatorname{Cont}$ means ordinary unrestricted recovery with the frozen PPO implementation, not an additional constraint that later fields remain equal. $N$ is the native continuation; $R$ the actor-restored continuation; $Q_D$ the matched onset-\emph{learner} reference; and $B_D$ the reciprocal impaired-actor continuation. $Q_D$ is neither the entire onset runner $P$ nor the clean removal runner $C^{+}$.

The archive's three-letter codes select actor, critic and Adam donors in that order. In \emph{pause cells}, C selects the onset donor $P$, not $C^{+}$. Thus N/R/$Q_D$/$B_D$ correspond to DDD/CDD/CCC/DCC under natural D context. The separate \texttt{fork/clean/donor/CCC} cell instead denotes the actual clean removal fork. These labels are an implementation convention, not interchangeable definitions of a clean reference. Appendix~\ref{app:compositions} lists all 11 cells and every preserved field.

\begin{figure}[p]
\centering
\resizebox{\linewidth}{!}{% One diagram, generated directly from the declared state-composition rules.
% Included inside a figure; source coordinates are centimetres.
\begin{tikzpicture}[x=1cm,y=1cm,>=Stealth,
  box/.style={draw=black!55,rounded corners=2pt,align=center,font=\small,inner sep=5pt},
  pbox/.style={box,fill=teal!10,draw=teal!75!black},
  dbox/.style={box,fill=red!7,draw=red!65!black},
  hbox/.style={box,dashed,draw=blue!65!black,fill=blue!3},
  cellP/.style={pbox,minimum width=2.1cm,minimum height=.62cm,font=\normalsize},
  cellD/.style={dbox,minimum width=2.1cm,minimum height=.62cm,font=\normalsize}]
\node[anchor=west,font=\bfseries] at (0,1.0) {A\quad Observed forks and a constructed reference};
\node[pbox,text width=2.7cm] (P) at (1.5,0) {$P$: shared onset\\rollout update 390};
\node[box,fill=gray!8,text width=3.0cm] (C) at (7.0,0) {$C^{+}$: clean removal\\update 440\\extra clean progress};
\node[box,fill=gray!8,text width=3.0cm] (trace) at (13.6,0) {Recorded clean trace\\$C(u)$, $u=390\ldots640$};
\node[dbox,text width=3.0cm] (D) at (7.0,-1.9) {$D$: impaired removal\\update 440\\natural context $x_D$};
\node[hbox,text width=3.2cm] (H) at (13.6,-1.9) {$H(k)=C(390+k)$\\historical benchmark\\not the matched $Q_D$};
\draw[->,thick] (P) -- node[above,font=\scriptsize]{50 clean updates} (C);
\draw[->,thick] (C) -- (trace);
\draw[->,thick,red!65!black] (P.south) to[out=-45,in=180] (D.west);
\node[font=\footnotesize,text=red!65!black,align=center,text width=3.0cm] at (3.7,-1.25) {50 reward-sign-\\corrupted updates};
\draw[->,dashed,blue!65!black,thick] (trace) -- node[right,font=\scriptsize]{shift the time index} (H);
\node[font=\footnotesize,align=left,anchor=west,text width=15.8cm] at (0,-3.15) {\textcolor{teal!70!black}{Repair uses $P$, not the extra-progress $C^{+}$.} $H$ reindexes the observed clean trace; it is not assumed to match a learner paused in the D context.};

\node[anchor=west,font=\bfseries] at (0,-4.1) {B\quad Initial fields of the executed interventions};
\node[font=\small] at (1.35,-4.85) {arm (archive code)};
\node[font=\small] at (4.3,-4.85) {actor parameters};
\node[font=\small] at (7.0,-4.85) {critic parameters};
\node[font=\small] at (9.7,-4.85) {Adam state};
\node[font=\small] at (13.6,-4.85) {other numerical runner fields};

\node[anchor=west,font=\small] at (0,-5.5) {$N$\quad(DDD)};
\node[cellD] at (4.3,-5.5) {$a_D$}; \node[cellD] at (7.0,-5.5) {$v_D$}; \node[cellD] at (9.7,-5.5) {$o_D$};
\node[anchor=west,font=\small] at (0,-6.3) {$R$\quad(CDD)};
\node[cellP] at (4.3,-6.3) {$a_P$}; \node[cellD] at (7.0,-6.3) {$v_D$}; \node[cellD] at (9.7,-6.3) {$o_D$};
\node[anchor=west,font=\small] at (0,-7.1) {$Q_D$\quad(CCC)};
\node[cellP] at (4.3,-7.1) {$a_P$}; \node[cellP] at (7.0,-7.1) {$v_P$}; \node[cellP] at (9.7,-7.1) {$o_P$};
\node[anchor=west,font=\small] at (0,-7.9) {$B_D$\quad(DCC)};
\node[cellD] at (4.3,-7.9) {$a_D$}; \node[cellP] at (7.0,-7.9) {$v_P$}; \node[cellP] at (9.7,-7.9) {$o_P$};
\node[box,fill=gray!6,text width=4.0cm,minimum height=3.2cm,font=\small] at (13.6,-6.7) {Same $x_D$ at removal:\\environment + observations,\\future PRNG key,\\rollout-update counter,\\\mbox{TrainState} step,\\other auxiliaries.\\\emph{Adam clock follows $o$.}};
\node[font=\footnotesize,anchor=west,text width=15.8cm] at (0,-9.0) {In the pause-cell archive code, C means an onset donor from $P$, \emph{not} the extra-progress clean fork $C^{+}$. $Q_D$ rolls back learner parameters/moments but retains the natural impaired context; it is not the full runner $P$.};

\node[anchor=west,font=\bfseries] at (0,-9.95) {C\quad Normal recovery after the splice: a total downstream effect};
\node[box,text width=3.8cm] (policy) at (2.2,-10.85) {Changed policy / actions};
\node[box,text width=4.1cm] (data) at (7.8,-10.85) {Changed experience,\\critic and optimizer updates};
\node[box,text width=3.6cm] (future) at (13.5,-10.85) {Future policy and\\evaluation returns};
\draw[->,thick] (policy) -- (data); \draw[->,thick] (data) -- (future);
\node[font=\small,align=center] at (8.0,-11.95) {Native cost: $A(N)-A(H)$\qquad Repair: $A(R)-A(N)$\\Reciprocal transfer: $A(B_D)-A(Q_D)$};
\node[font=\footnotesize,align=center,text width=15.8cm] at (8.0,-12.9) {Solid arrows: observed/executed training. Dashed arrow: a reference readout. Matching holds \emph{at removal}; future states and data are allowed to diverge. No arrow identifies uniquely stored impairment information.};
\end{tikzpicture}
}
\caption{\textbf{Donors, executed swaps and reference trajectories in one diagram.} A: $P$ precedes impairment; the extra-progress $C^{+}$ actor never supplies the primary repair. The historical $H$ trace is obtained by reindexing observations from the clean branch. B: all four primary continuations initially share $x_D$; colored fields show the P/D source of the actor, critic and Adam state. The code letter C denotes P only in these pause cells. C: actor replacement can change actions, experience and subsequent learner updates, so its effect includes these downstream paths. Solid arrows mark actual training/continuations; the dashed arrow marks a reference operation. Whole-runner freezing and learner rollback in the D world are different counterfactuals even when their mean returns are close.}
\label{m:counterfactuals}
\end{figure}

\subsection{A historical reference is not a unique literal pause}\label{m:pausemeaning}
Let $C_i(u)$ denote the observed unimpaired evaluation trajectory at training update $u$. Our historical reference is the specified time shift $H_i(k)=C_i(w_0+k)$, read at the same recovery index $k$ as the impaired trajectory. At removal, $H_i(0)$ is the pre-impairment clean evaluation. At recovery $k=200$, $H_i(200)$ uses clean update 590, whereas the native learner has reached update 640. This construction charges the impairment for the 50 updates it occupies rather than crediting the comparator with those extra updates.

There is a literal pause implementation that \emph{can} realize this reference: freeze the \emph{entire} runner P, including environment, random stream and all clocks, allow no external process to change it, and resume unrestricted training with the same deterministic, stationary program. Ignoring elapsed wall time, this produces the same clean continuation from P. This interpretation requires those specifications; ``pause'' alone does not determine them.

A paused learner resumed in the impaired learner's removal-time world is different. An environment might continue evolving during the pause, its data might be retained or discarded, and training/RNG/optimizer clocks might be frozen or advanced. Each choice defines another state and another counterfactual. $Q_D$ is one explicitly executed learner-rollback intervention in $x_D$, not a claim that an actual pause would naturally produce $x_D$. Indeed, $x_D$ was generated by the impaired learner's actions. No experiment equates these counterfactuals by naming both of them a pause.

Consequently, ``native recovery cost'' always means \emph{native cost relative to the declared historical H benchmark}. It is not an intrinsic scalar attached to the impaired state independently of a reference. A small difference between the mean returns of $Q_D$ and H can support a chosen practical margin in this assay; it does not prove identical states, identical trajectories or future-context neutrality. We therefore report repaired-minus-H and repaired-minus-$Q_D$ residuals separately, including the failed latter margin.

\subsection{A total downstream intervention, not uniquely stored impairment}\label{m:total_effect}
Let $\Phi^k(S)$ be the complete training continuation for $k$ rollout updates from state S, including policy-dependent environment interaction and learner updates. The actor-restoration contrast compares evaluation of $\Phi^k\{\mathcal S(a_P,v_D,o_D;x_D)\}$ with evaluation of $\Phi^k\{\mathcal S(a_D,v_D,o_D;x_D)\}$. The root difference is actor parameters; the resulting effect is \emph{total}, not restricted to the initial actor's direct contribution to a fixed dataset or fixed future state sequence.

The paths through changed action choices, visited states, rewards, critic estimates and optimizer updates are allowed to operate. After the first subsequent action, the experience used by the two learners need not coincide, even with common initial random keys. Our intervention does not block those mediators and does not estimate a controlled direct effect holding them fixed. It therefore establishes a useful causal handle under a specified background, not a natural decomposition of encoded information.

Nor does a successful actor transplant show that critic/optimizer state is uninvolved. Those fields may interact with the transplanted actor, and their conditional effects can change with the chosen donor or context. The smaller negative critic effect under common fresh Adam and the incomplete restoration residual are retained rather than absorbed into an ``all damage is in the actor'' conclusion. A larger reference-dependent repair ratio would still not identify a unique internal representation of impairment.
